% GENERATED FILE. Edit canonical lessons, then run npm run generate:books.
% canonical-source-hash: fnv1a64:e86b035ff71f4e65
% canonical-lessons: LA-C140-pecten, LA-C140-forfex, LA-C140-acus, LA-C140-filum, LA-C140-lana, LA-R140-first-pass-food-and-nature, LA-R140-second-pass-animals-and-the-body

\chapter{Comb, Scissors, Needle, Thread, Wool}
\label{ch:la-comb-scissors-needle-thread-wool}
\hlchaptermodality{140}

\begin{chapteropening}
\textbf{By the end of this chapter:} \emph{I can say a comb, scissors, a needle, a thread and wool.}

Complete the last lesson of chapter 140: I can say a comb, scissors, a needle, a thread and wool.
\end{chapteropening}

\section[pecten, pectinis]{pecten, pectinis — a comb}

\label{lesson:LA-C140-pecten}



\begin{quote}
Before the new one: say the Latin for a torch, then the Latin for a lantern.
\end{quote}

\subsection*{You'll want to know: pecten, pectinis}
\textbf{pecten, pectinis} — \textquotedblleft{}a comb\textquotedblright{}.

\begin{cousinweb}
\textbf{pecten} is built on \emph{pectō}, \textquotedblleft{}to comb\textquotedblright{} (Lewis \& Short), and is usually traced to Proto-Indo-European \emph{*pek-}, \textquotedblleft{}to pluck wool or hair\textquotedblright{} (EtymDB 2.1).
\end{cousinweb}

\begin{grammarlens}[title={What you've built}]
One more word for this chapter.
\end{grammarlens}

\subsection*{Your turn}
\begin{itemize}
\raggedright
  \item \emph{Say it:} \emph{pecten}
  \item \emph{Say it:} \emph{pecten}, once more
  \item \emph{Say it:} say \emph{pecten}
  \item \emph{Recall:} say the Latin for a torch, then the Latin for a lantern, then say \emph{pecten} again
\end{itemize}

\begin{tcolorbox}[breakable,colback=teal!4,colframe=teal!35!black,title={Before you move on}]
What does pecten, pectinis mean? (\textquotedblleft{}A comb\textquotedblright{}.) Say it once more.
\end{tcolorbox}
\section[forfex, forficis]{forfex, forficis — scissors}

\label{lesson:LA-C140-forfex}



\begin{quote}
Before the new one: say the Latin for a lantern, then the Latin for a comb.
\end{quote}

\subsection*{You'll want to know: forfex, forficis}
\textbf{forfex, forficis} — \textquotedblleft{}scissors\textquotedblright{}.

\begin{cousinweb}
\textbf{forfex} has no sure origin; it is often linked with \emph{forceps}, \textquotedblleft{}tongs\textquotedblright{}, but the link is uncertain.
\end{cousinweb}

\begin{grammarlens}[title={What you've built}]
One more word for this chapter.
\end{grammarlens}

\subsection*{Your turn}
\begin{itemize}
\raggedright
  \item \emph{Say it:} \emph{forfex}
  \item \emph{Say it:} \emph{forfex}, once more
  \item \emph{Say it:} say \emph{forfex}
  \item \emph{Recall:} say the Latin for a lantern, then the Latin for a comb, then say \emph{forfex} again
\end{itemize}

\begin{tcolorbox}[breakable,colback=teal!4,colframe=teal!35!black,title={Before you move on}]
What does forfex, forficis mean? (\textquotedblleft{}Scissors\textquotedblright{}.) Say it once more.
\end{tcolorbox}
\section[acus, acūs]{acus, acūs — a needle}

\label{lesson:LA-C140-acus}



\begin{quote}
Before the new one: say the Latin for a comb, then the Latin for scissors.
\end{quote}

\subsection*{You'll want to know: acus, acūs}
\textbf{acus, acūs} — \textquotedblleft{}a needle\textquotedblright{}.

\begin{cousinweb}
\textbf{acus} is inherited, usually traced to Proto-Indo-European \emph{*h2ek-}, \textquotedblleft{}sharp\textquotedblright{} (EtymDB 2.1), the root of \emph{ācer}, \textquotedblleft{}sharp\textquotedblright{} (Lewis \& Short).
\end{cousinweb}

\begin{grammarlens}[title={What you've built}]
One more word for this chapter.
\end{grammarlens}

\subsection*{Your turn}
\begin{itemize}
\raggedright
  \item \emph{Say it:} \emph{acus}
  \item \emph{Say it:} \emph{acus}, once more
  \item \emph{Say it:} say \emph{acus}
  \item \emph{Recall:} say the Latin for a comb, then the Latin for scissors, then say \emph{acus} again
\end{itemize}

\begin{tcolorbox}[breakable,colback=teal!4,colframe=teal!35!black,title={Before you move on}]
What does acus, acūs mean? (\textquotedblleft{}A needle\textquotedblright{}.) Say it once more.
\end{tcolorbox}
\section[fīlum, fīlī]{fīlum, fīlī — a thread}

\label{lesson:LA-C140-filum}



\begin{quote}
Before the new one: say the Latin for scissors, then the Latin for a needle.
\end{quote}

\subsection*{You'll want to know: fīlum, fīlī}
\textbf{fīlum, fīlī} — \textquotedblleft{}a thread\textquotedblright{}.

\begin{cousinweb}
\textbf{fīlum} is inherited, usually traced to Proto-Indo-European \emph{*gwhiHslo-}, \textquotedblleft{}thread\textquotedblright{} (EtymDB 2.1). English \emph{filament} and \emph{file}, a row, come from it.
\end{cousinweb}

\begin{grammarlens}[title={What you've built}]
One more word for this chapter.
\end{grammarlens}

\subsection*{Your turn}
\begin{itemize}
\raggedright
  \item \emph{Say it:} \emph{fīlum}
  \item \emph{Say it:} \emph{fīlum}, once more
  \item \emph{Say it:} say \emph{fīlum}
  \item \emph{Recall:} say the Latin for scissors, then the Latin for a needle, then say \emph{fīlum} again
\end{itemize}

\begin{tcolorbox}[breakable,colback=teal!4,colframe=teal!35!black,title={Before you move on}]
What does fīlum, fīlī mean? (\textquotedblleft{}A thread\textquotedblright{}.) Say it once more.
\end{tcolorbox}
\section[lāna, lānae]{lāna, lānae — wool}

\label{lesson:LA-C140-lana}



\begin{quote}
Before the new one: say the Latin for a needle, then the Latin for a thread.
\end{quote}

\subsection*{You'll want to know: lāna, lānae}
\textbf{lāna, lānae} — \textquotedblleft{}wool\textquotedblright{}.

\begin{cousinweb}
\textbf{lāna} is inherited, usually traced to Proto-Indo-European \emph{*h2wlh1neh2}, \textquotedblleft{}wool\textquotedblright{} (EtymDB 2.1); Latin lost the first w. English \emph{wool} is a cousin.
\end{cousinweb}

\begin{grammarlens}[title={What you've built}]
One more word for this chapter.
\end{grammarlens}

\subsection*{Your turn}
\begin{itemize}
\raggedright
  \item \emph{Say it:} \emph{lāna}
  \item \emph{Say it:} \emph{lāna}, once more
  \item \emph{Say it:} say \emph{lāna}
  \item \emph{Recall:} say the Latin for a needle, then the Latin for a thread, then say \emph{lāna} again
\end{itemize}

\begin{tcolorbox}[breakable,colback=teal!4,colframe=teal!35!black,title={Before you move on}]
What does lāna, lānae mean? (\textquotedblleft{}Wool\textquotedblright{}.) Say it once more.
\end{tcolorbox}
\section[(dialogue)]{Food and nature — a first pass}

\label{lesson:LA-R140-first-pass-food-and-nature}



\begin{quote}
Say the words for a thread and wool.
\end{quote}

\subsection*{Your turn}
Say these aloud:

\begin{itemize}
\raggedright
  \item fog, an eagle, a dove, a raven, a duck
  \item a goose, a bee, a fly, an ant, a butterfly
  \item a spider, a frog, a hare, a fox, a bull
  \item a cow, a mule, a camel, a monkey, a dolphin
  \item a whale, a cockerel, a lip, the chin, the back
\end{itemize}

\begin{tcolorbox}[breakable,colback=teal!4,colframe=teal!35!black,title={Before you move on}]
Fog? (\textbf{nebula, nebulae}.) An eagle? (\textbf{aquila, aquilae}.) A dove? (\textbf{columba, columbae}.) A raven? (\textbf{corvus, corvī}.) A duck? (\textbf{anas, anatis}.) A goose? (\textbf{ānser, ānseris}.) A bee? (\textbf{apis, apis}.) A fly? (\textbf{musca, muscae}.) An ant? (\textbf{formīca, formīcae}.) A butterfly? (\textbf{pāpiliō, pāpiliōnis}.) A spider? (\textbf{arānea, arāneae}.) A frog? (\textbf{rāna, rānae}.) A hare? (\textbf{lepus, leporis}.) A fox? (\textbf{vulpēs, vulpis}.) A bull? (\textbf{taurus, taurī}.) A cow? (\textbf{vacca, vaccae}.) A mule? (\textbf{mūlus, mūlī}.) A camel? (\textbf{camēlus, camēlī}.) A monkey? (\textbf{sīmia, sīmiae}.) A dolphin? (\textbf{delphīnus, delphīnī}.) A whale? (\textbf{balaena, balaenae}.) A cockerel? (\textbf{gallus, gallī}.) A lip? (\textbf{labrum, labrī}.) The chin? (\textbf{mentum, mentī}.) The back? (\textbf{dorsum, dorsī}.)
\end{tcolorbox}
\section[(dialogue)]{Animals and the body — a second pass}

\label{lesson:LA-R140-second-pass-animals-and-the-body}



\begin{quote}
Say the words for an elbow and a thumb.
\end{quote}

\subsection*{Your turn}
Say these aloud:

\begin{itemize}
\raggedright
  \item an elbow, a thumb, the skin, a bone, the stomach
  \item a craftsman, a baker, a shoemaker, a barber, a shepherd
  \item a fisherman, a slave, a mistress, a friend, a cupboard
  \item a chair with a back, a bench, a carpet, a torch, a lantern
  \item a comb, scissors, a needle, a thread, wool
\end{itemize}

\begin{tcolorbox}[breakable,colback=teal!4,colframe=teal!35!black,title={Before you move on}]
An elbow? (\textbf{cubitum, cubitī}.) A thumb? (\textbf{pollex, pollicis}.) The skin? (\textbf{cutis, cutis}.) A bone? (\textbf{os, ossis}.) The stomach? (\textbf{stomachus, stomachī}.) A craftsman? (\textbf{faber, fabrī}.) A baker? (\textbf{pistor, pistōris}.) A shoemaker? (\textbf{sūtor, sūtōris}.) A barber? (\textbf{tōnsor, tōnsōris}.) A shepherd? (\textbf{pāstor, pāstōris}.) A fisherman? (\textbf{piscātor, piscātōris}.) A slave? (\textbf{servus, servī}.) A mistress? (\textbf{domina, dominae}.) A friend? (\textbf{amīca, amīcae}.) A cupboard? (\textbf{armārium, armāriī}.) A chair with a back? (\textbf{cathedra, cathedrae}.) A bench? (\textbf{scamnum, scamnī}.) A carpet? (\textbf{tapēte, tapētis}.) A torch? (\textbf{fax, facis}.) A lantern? (\textbf{lanterna, lanternae}.) A comb? (\textbf{pecten, pectinis}.) Scissors? (\textbf{forfex, forficis}.) A needle? (\textbf{acus, acūs}.) A thread? (\textbf{fīlum, fīlī}.) Wool? (\textbf{lāna, lānae}.)
\end{tcolorbox}
